%%
%% Author: shoupu_wan
%% 10/10/18
%%

\chapter{Case Study: $N$ Queens}\label{ch:n-queens}

\lstinputlisting[]{loop/word_breaker.py}
\lstinputlisting[]{loop/word_breaker_other_imp.py}


\section{Solution}
Our solution is backtracking.
We are going to iterate through the solutions in digital-clock order.
We place the first queen at the first possible position.
If current placement is feasible, we keep exploring the current solution space.
If we encountered dead-end, we backtrack to previous working solution by removing the last queen.
If our first queen is moved off board (at position 0), then we are done with iteration.

This is similar to the `next permutation' problem.


\lstinputlisting[label={lst:OldNQueensSolution},
caption={Count the number of solutions for the `$N$-queens' problem.}]
{case-study/n-queens/NQueens_old.java}

After refactor to better SRP and shrunk parameter list for functions \code{place()} and
\code{check()}.

\lstinputlisting[label={lst:NQueensSolution},
caption={Refactored code based on ~\autoref{lst:OldNQueensSolution} for the `$N$-queens' problem.}]
{case-study/n-queens/NQueens.java}

If one want to list solutions instead of counting them, the changes I made on top of the
~\autoref{lst:NQueensSolution} are shown in ~\autoref{lst:NQueensSolutionListSolutions}.

\lstinputlisting[label={lst:NQueensSolutionListSolutions},
caption={List solutions for the `$N$-queens' problem.}]
{case-study/n-queens/NQueensListSolutions.java}
